\documentclass[10pt,a4paper]{amsart}
\usepackage{amsmath,amssymb,amsthm,graphicx}
\usepackage{verbatim}
\newenvironment{reference}{\comment}{\endcomment}
%\newenvironment{reference}{}{}
\newenvironment{slogan}{\comment}{\endcomment}
\newenvironment{history}{\comment}{\endcomment}
% For commutative diagrams we use Xy-pic
\usepackage[all]{xy}

% We use 2cell for 2-commutative diagrams.
\xyoption{2cell}
% Theorem environments.
%
\theoremstyle{plain}
\newtheorem{theorem}[subsection]{Theorem}
\newtheorem{proposition}[subsection]{Proposition}
\newtheorem{lemma}[subsection]{Lemma}

\theoremstyle{definition}
\newtheorem{definition}[subsection]{Definition}
\newtheorem{example}[subsection]{Example}
\newtheorem{exercise}[subsection]{Exercise}
\newtheorem{situation}[subsection]{Situation}

\theoremstyle{remark}
\newtheorem{remark}[subsection]{Remark}
\newtheorem{remarks}[subsection]{Remarks}

\numberwithin{equation}{subsection}

% Macros
%
\def\lim{\mathop{\mathrm{lim}}\nolimits}
\def\colim{\mathop{\mathrm{colim}}\nolimits}
\def\Spec{\mathop{\mathrm{Spec}}}
\def\Hom{\mathop{\mathrm{Hom}}\nolimits}
\def\Ext{\mathop{\mathrm{Ext}}\nolimits}
\def\SheafHom{\mathop{\mathcal{H}\!\mathit{om}}\nolimits}
\def\SheafExt{\mathop{\mathcal{E}\!\mathit{xt}}\nolimits}
\def\Sch{\mathit{Sch}}
\def\Mor{\mathop{\mathrm{Mor}}\nolimits}
\def\Ob{\mathop{\mathrm{Ob}}\nolimits}
\def\Sh{\mathop{\mathit{Sh}}\nolimits}
\def\NL{\mathop{N\!L}\nolimits}
\def\CH{\mathop{\mathrm{CH}}\nolimits}
\def\proetale{{pro\text{-}\acute{e}tale}}
\def\etale{{\acute{e}tale}}
\def\QCoh{\mathit{QCoh}}
\def\Ker{\mathop{\mathrm{Ker}}}
\def\Im{\mathop{\mathrm{Im}}}
\def\Coker{\mathop{\mathrm{Coker}}}
\def\Coim{\mathop{\mathrm{Coim}}}

% Boxtimes
%
\DeclareMathSymbol{\boxtimes}{\mathbin}{AMSa}{"02}

%
% Macros for moduli stacks/spaces
%
\def\QCohstack{\mathcal{QC}\!\mathit{oh}}
\def\Cohstack{\mathcal{C}\!\mathit{oh}}
\def\Spacesstack{\mathcal{S}\!\mathit{paces}}
\def\Quotfunctor{\mathrm{Quot}}
\def\Hilbfunctor{\mathrm{Hilb}}
\def\Curvesstack{\mathcal{C}\!\mathit{urves}}
\def\Polarizedstack{\mathcal{P}\!\mathit{olarized}}
\def\Complexesstack{\mathcal{C}\!\mathit{omplexes}}
% \Pic is the operator that assigns to X its picard group, usage \Pic(X)
% \Picardstack_{X/B} denotes the Picard stack of X over B
% \Picardfunctor_{X/B} denotes the Picard functor of X over B
\def\Pic{\mathop{\mathrm{Pic}}\nolimits}
\def\Picardstack{\mathcal{P}\!\mathit{ic}}
\def\Picardfunctor{\mathrm{Pic}}
\def\Deformationcategory{\mathcal{D}\!\mathit{ef}}
\usepackage{cleveref}
\crefname{originalreference}{original reference}{original references}
\Crefname{originalreference}{Original reference}{Original references}
\makeatletter
\newlabel{criteria-lemma-flat-quotient-flat-presentation}{{criteria-lemma-flat-quotient-flat-presentation}{source}{}{sourcepointer-1}{}}
\newlabel{criteria-lemma-flat-quotient-flat-presentation@cref}{{[originalreference][0][]criteria-lemma-flat-quotient-flat-presentation}{[1][0][]source}{}{}{}}
\makeatother
\crefname{originalreference}{original reference}{original references}
\Crefname{originalreference}{Original reference}{Original references}
\makeatletter
\newlabel{criteria-theorem-flat-groupoid-gives-algebraic-stack}{{criteria-theorem-flat-groupoid-gives-algebraic-stack}{source}{}{sourcepointer-2}{}}
\newlabel{criteria-theorem-flat-groupoid-gives-algebraic-stack@cref}{{[originalreference][0][]criteria-theorem-flat-groupoid-gives-algebraic-stack}{[1][0][]source}{}{}{}}
\makeatother
\crefname{originalreference}{original reference}{original references}
\Crefname{originalreference}{Original reference}{Original references}
\makeatletter
\newlabel{spaces-descent-lemma-smooth-permanence}{{spaces-descent-lemma-smooth-permanence}{source}{}{sourcepointer-3}{}}
\newlabel{spaces-descent-lemma-smooth-permanence@cref}{{[originalreference][0][]spaces-descent-lemma-smooth-permanence}{[1][0][]source}{}{}{}}
\makeatother
\title{A quotient of a smooth space by a flat finitely presented group is smooth over the base}
\author{The Stacks Project Authors}
\date{stacks-project:stacks-morphisms.tex:local-snapshot-20261012; GNU FDL 1.2 or later}
\begin{document}
\maketitle
\noindent\textit{Editorial note.} All delivered mathematical text is the author's own, copied from the pinned Stacks Project chapter source sources/stacks-algebraic-stacks-morphisms/pinned-original-source.tex; no statement or proof step was rewritten, reordered or removed.\par\medskip
\noindent\textit{Editorial note.} Preamble provenance: the Stacks Project keeps one single shared preamble for every chapter; this card takes it directly from the staged sources/stacks-algebraic-stacks-morphisms/preamble.tex as three verbatim spans (lines 8-12, 14-18 and 154-217), so no derived file was written and nothing was added or redefined. The wrapper class and the base packages are supplied by the assembler.\par\medskip
\noindent\textit{Editorial note.} Bibliography provenance: neither the counted statement nor its proof contains a citation (this lane rescans the whole span before writing the card), so the card carries no bibliography and no entry was reconstructed.\par\medskip
\noindent\textit{Editorial note.} Selection provenance: this card was opened under the file-based control contract (CONTROL.md), with the mechanical pre-card gate run on its statement and proof interval before the card was written.\par\medskip
\noindent\textit{Allowed standard prerequisites (editorial).}\begin{itemize}
\item The theorem that a flat groupoid in algebraic spaces gives an algebraic stack. Invoked at the start of the counted proof (source lines 7709-7711) as a cross-chapter pointer. Makes the quotient stack algebraic.
\item The statement that the covering of the quotient is flat and locally of finite presentation. Invoked inside the counted proof (source lines 7712-7718) as a cross-chapter pointer. Supplies the base change used to apply descent for smoothness.
\item The descent statement for smoothness along a flat finitely presented surjection. Invoked at the end of the counted proof (source lines 7719-7724) as a cross-chapter pointer. Concludes smoothness of the quotient over the base.
\end{itemize}\medskip
\ifdefined\hypertarget\hypertarget{sourcepointer-1}{}\fi
\noindent\textit{Source reference.} Original reference criteria-lemma-flat-quotient-flat-presentation: sources/stacks-algebraic-stacks-morphisms/pinned-original-source.tex, source lines 3877--3877. Navigation pointer only: the counted text refers to the Stacks item criteria-lemma-flat-quotient-flat-presentation, which belongs to another chapter of the project and is not part of this card. The pointer carries the target's own label as its number because this chapter's staged snapshot does not contain the other chapter's tag table; its content is not reproduced here.\par\medskip
\ifdefined\hypertarget\hypertarget{sourcepointer-2}{}\fi
\noindent\textit{Source reference.} Original reference criteria-theorem-flat-groupoid-gives-algebraic-stack: sources/stacks-algebraic-stacks-morphisms/pinned-original-source.tex, source lines 3850--3850. Navigation pointer only: the counted text refers to the Stacks item criteria-theorem-flat-groupoid-gives-algebraic-stack, which belongs to another chapter of the project and is not part of this card. The pointer carries the target's own label as its number because this chapter's staged snapshot does not contain the other chapter's tag table; its content is not reproduced here.\par\medskip
\ifdefined\hypertarget\hypertarget{sourcepointer-3}{}\fi
\noindent\textit{Source reference.} Original reference spaces-descent-lemma-smooth-permanence: sources/stacks-algebraic-stacks-morphisms/pinned-original-source.tex, source lines 7724--7724. Navigation pointer only: the counted text refers to the Stacks item spaces-descent-lemma-smooth-permanence, which belongs to another chapter of the project and is not part of this card. The pointer carries the target's own label as its number because this chapter's staged snapshot does not contain the other chapter's tag table; its content is not reproduced here.\par\medskip
\begin{lemma}
\label{lemma-smooth-quotient-stack}
Let $X \to Y$ be a smooth morphism of algebraic spaces.
Let $G$ be a group algebraic space over $Y$ which is flat
and locally of finite presentation over $Y$. Let $G$ act on $X$ over $Y$.
Then the quotient stack $[X/G]$ is smooth over $Y$.
\end{lemma}

\begin{proof}
The quotient $[X/G]$ is an algebraic stack by
Criteria for Representability, Theorem
\ref{criteria-theorem-flat-groupoid-gives-algebraic-stack}.
The smoothness of $[X/G]$ over $Y$ follows from the fact that smoothness
descends under fppf coverings:
Choose a surjective smooth morphism $U \to [X/G]$ where $U$ is a scheme.
Smoothness of $[X/G]$ over $Y$ is equivalent to smoothness of $U$ over $Y$.
Observe that $U \times_{[X/G]} X$ is smooth over $X$ and hence smooth
over $Y$ (because compositions of smooth morphisms are smooth).
On the other hand, $U \times_{[X/G]} X \to U$ is locally of
finite presentation, flat, and surjective (because it is
the base change of $X \to [X/G]$ which has those properties
for example by Criteria for Representability, Lemma
\ref{criteria-lemma-flat-quotient-flat-presentation}).
Therefore we may apply Descent on Spaces,
Lemma \ref{spaces-descent-lemma-smooth-permanence}.
\end{proof}

\end{document}
